\scanpage{8}
tel que le PL (ou d'autres meilleurs, à la \unsure{Hironaka} p.ex.), le cadre topologique pur ne s'y prêtant pas. Il semble néanmoins que dans le cas des \unsure{n-cartes|cartes} loc. finies \uline{régulières}, les difficultés disparaissent \deleted{\unsure{topologiques}}
\margin{à vérifier avec soin ....}

10) Les gpes apparaissant comme groupes d'automorphismes de ``cartes régulières \deleted{\unsure{(du pt de v\textsuperscript{u})}} \uline{simplement connexes} (loc finies ?), \deleted{\unsure{orientées}} \inserted{\unsure{orientées}} ou pas, sont tous des groupes de Coxeter. Ceux parmi eux qui sont finis sont bien connus (cf livre Bourbaki) -- il semble qu'ils correspondent tous à des sphères -- avec la décomposition cellulaire ``simpliciale'' évidente, liée à la géométrie des gpes de Coxeter.
\margin{id}

D'après Deligne, toutes les ``cartes régulières (finies ou non) \inserted{simplement connexes} devraient

\mspage{7}
